% ===== PDF p.026 (原书 p.18) =====
\chapter{Gromov--Witten 不变量入门}
\chapterauthor{深谷賢治}{[京都大学大学院理学研究科]}

\section{引言}

Gromov--Witten 不变量（Gromov--Witten invariant）正如其名，是由 Gromov 和 Witten 从彼此不同的动机出发各自独立发现的。

Gromov 的动机在辛几何。研究大范围的辛几何，一个基本手段就是研究从黎曼面到近复流形的全纯映射。Gromov 的这一想法后来由 Ruan 等人加以深化，确立了这样一种不变量：数出从黎曼面到辛流形的拟全纯映射（pseudoholomorphic map）的个数。这就是今天被称为 Gromov--Witten 不变量的东西。

另一方面，Witten 则从带有超对称的拓扑 sigma 模型的配分函数出发，引入了同一个不变量。

从这种物理的观点看，举例来说，Gromov--Witten 不变量并不只是（近）复流形的不变量，它更是辛流形的不变量。不过，我觉得它在辛几何上的这个应用，最根本的一点反而没有被理解。

另一方面，把 Gromov--Witten 不变量当作拓扑 sigma 模型的配分函数来看，最大的好处在于：它的许多性质都能从物理的考察推导出来。典型的例子是结合律。也就是说，由 Gromov--Witten 不变量得到的上同调环的形变（称为量子上同调环）仍是环，尤其形变之后乘法仍满足结合律，这一点很重要。把这一点一般化，便得到这样一个微分方程：由 Gromov--Witten 不变量构造的母函数满足它。这就是 WDVV 方程\footnote{取 Witten、Dijkgraaf、Verlinde、Verlinde 四人的首字母命名。}。此外，本书的主题——镜像对称——中，Gromov--Witten 不变量之所以作为一方主角登场，
% ===== PDF p.027 (原书 p.19) =====
其理由也只有在把 Gromov--Witten 不变量看作拓扑 sigma 模型的配分函数之后，才能得到理解。此外，第 1 章里出现的 Virasoro 猜想等也是这样的例子。

关于 Gromov--Witten 不变量，笔者已经写过几篇解说（[1] 的 135--174 页、[2]、[3] 的第 6 章、[4] 的第 V 章等）。所以本章讲的内容比它们浅显得多：先复习什么是相交数、什么是同调，再讲一个最简单的例子，看看由 Gromov--Witten 不变量导出的量子上同调与普通上同调有什么不同。那就是复射影平面的情形——它变成什么，又对应着怎样的几何事实，等等。

\section{相交数}

Gromov--Witten 不变量可以看作相交理论的一种。先复习相交理论。

\begin{theorem*}{2.1}
$n$ 次方程
\begin{equation}
x^n + a_1 x^{n-1} + \cdots + a_{n-1} x + a_n = 0 \tag{1}
\end{equation}
的根，把重根也算上、并在复数范围内考虑，恰好有 $n$ 个。
\end{theorem*}

这个定理是高中就学过的。方程（1）可以改写成联立方程组
\begin{equation}
y = x^n + a_1 x^{n-1} + \cdots + a_{n-1} x + a_n, \qquad y = 0 \tag{2}
\end{equation}
把它一般化，我们来考虑
\begin{equation}
\begin{cases}
P(x,y) = 0 \\
Q(x,y) = 0
\end{cases} \tag{3}
\end{equation}
这里 $P,Q$ 分别是 $n,m$ 次的（二元）多项式。

% ===== PDF p.028 (原书 p.20) =====
\begin{theorem*}{2.2（Bézout 定理）}
方程（3）的根的个数恰好是 $mn$ 个。不过要把重根算上，而且“无穷远处”的根也要数进去。
\end{theorem*}

那么，上面说的 Bézout 定理意味着：

\begin{center}
\textbf{方程（3）的根的个数是“拓扑的”量}
\end{center}

下面来说明这件事。考虑两个集合
\[
\cP = \{(x,y) \mid P(x,y) = 0\}, \qquad
\cQ = \{(x,y) \mid Q(x,y) = 0\}.
\]
（3）的根的个数就相当于
\begin{equation}
\#(\cP \cap \cQ) \tag{4}
\end{equation}
（$\# A$ 表示集合 $A$ 中元素的个数。）

像（4）这样的数称为“相交数”（intersection number）。Bézout 定理主张的是：相交数只由表示图形的方程的次数决定。次数相同的两个图形，通过连续地改变系数，可以从一方连续变形到另一方。也就是说：

\begin{center}
\textbf{相交数在图形的连续变形下不变}
\end{center}

这就是“相交数是一个拓扑的量”这句话的意思。

这里说明一下“无穷远处的相交也要考虑”这句话的含义。对常数 $a$ 取各种值，来考虑联立方程组
\begin{equation}
\begin{cases}
2x + 3y = 0 \\
6x + ay = 3
\end{cases} \tag{5}
\end{equation}
只要 $a \neq 9$，（5）就恰好有一个解。（5）是（4）的特例，$n = m = 1$。所以这里的 1 就是 $1 \times 1$，与 Bézout 定理对得上。

问题出在 $a = 9$ 的情形，这时（5）没有解。这样不方便，所以当 $a = 9$ 时，我们就认为“无穷远处有解”。

所谓射影空间（现在考虑的是复射影平面），就是设想每一直线方向各有一个“无穷远点”，再把这些无穷远点添到平面上去。在现在的例子里，$2x + 3y = 0$ 的方向（写成向量就是 $(3,-2)$ 的方向）所对应的无穷远点，就认为 $\cP$ 与 $\cQ$ 相交。

% [Tutor Note: 原文此处写“（5）是（4）的特殊情形”，但（4）是相交数 $\#(P\cap Q)$，（3）才是 $P,Q$ 的一般方程组；按文意疑应为（3）。此处照原文录为（4）。]

% ===== PDF p.029 (原书 p.21) =====
\section{同调}

前一节从代数几何的角度，也就是考虑多项式解的集合，讲述了相交数。若要讲得更拓扑一些，就要使用“同调”（homology）这个说法。

对空间 $X$，考虑 $X$ 中所含 $d$ 维图形的全体所成的集合 $\mathcal{X}(d)$。给定 $\mathcal{X}(d)$ 的两个元素 $\cP$ 与 $\cQ$ 之后，我们想考虑它们的相交数；可是 $\mathcal{X}(d)$ 作为集合太大了，在那里考虑也弄不清楚。

于是，当 $\mathcal{X}(d)$ 的元素 $\cP$ 与另一个元素 $\cP'$“可以通过连续变形互移”\footnote{“通过连续变形互移”这种说法，拿来解说同调群是相当不准确的。正确的说明不是“连续变形可互移”，而是“同调”（homologous），这里从略。}时，就把 $\cP$ 与 $\cP'$ 视为同一个东西。这样一来，就得到了比 $\mathcal{X}(d)$ 小得多、也容易处理的集合。这大体上就是 $X$ 的 $d$ 次同调群。

例如，在复射影平面中考虑多项式 $P(x,y)$ 的根全体 $\cP$：若把可以连续变形互移的视为相同，那么剩下的就只由 $P$ 的次数 $k$ 决定。把这个元素记作 $k[\CP^1]$\footnote{$\CP^1$ 是（某个）1 次方程的根全体所表示的集合，也就是直线。}。

这里写成 $k[\CP^1]$，看起来像是把 $[\CP^1]$ 乘了 $k$ 倍，但确实可以这样考虑。例如 $k=2$，也就是 $P$ 是 2 次式的情形：
\[
P_a(x,y) = x^2 + axy + 4y^2 - 1
\]
让 $a$ 取各种值来考虑。$a=0$ 时，$P_a(x,y)$ 的根全体 $\cP_a$ 是椭圆。另一方面，让 $a$ 变动，当 $a=-4$ 时
\[
P_{-4}(x,y) = (x-2y+1)(x-2y-1),
\]
所以 $P_{-4}(x,y)$ 是直线 $\mathcal{L}_1 : x-2y+1=0$ 与直线 $\mathcal{L}_2 : x-2y-1=0$ 的并集。$\mathcal{L}_1$ 可以连续地移到 $\mathcal{L}_2$，所以视为同一个东西。结果：

% ===== PDF p.030 (原书 p.22) =====
\[
(\text{2 次式所表示的图形}) = 2 \times (\text{1 次式所表示的图形}),
\]
就得到这样的等式。（这里，集合的并在不知不觉之间变成了加法，不过就当它没问题吧。）更一般地，
\[
(\text{$k$ 次式所表示的图形}) = k \times (\text{1 次式所表示的图形}),
\]
这个等式成立。换句话说，“在复射影平面中，对每个整数 $k$，都有与（$k$ 次式所表示的图形）相对应的同调的元素。”把这个“……”中的话再说一遍，就是“\textbf{复射影平面的 2 次同调群是整数全体 $\Z$}”\footnote{多项式的次数不能是负数，所以 $k$ 不应是整数而应是自然数——这话固然不错。但在同调里，我们还要考虑不能由多项式表示的图形；而且相交数本来（一般地）还要考虑符号。由于这些情况，“复射影平面的 1 次同调”不是自然数全体，而是整数全体。}。空间 $X$ 的 $d$ 次同调群记作 $H_d(X)$。复射影平面 $\CP^2$ 的同调是
% [Tutor Note: 本页黑体断言作「2 次同调群」，而脚注 4 中同一对象却印作「1 次」。$u_2$ 表示直线，对应的同调是实 2 维的，脚注的「1 次」疑为笔误；两处均照原文录出。]
\begin{equation}
H_d(\CP^2) =
\begin{cases}
\Z, & d = 0,2,4, \\
0, & \text{其他情形}
\end{cases}
\tag{6}
\end{equation}

（6）中 $d=2$ 的情形已经说明过了。$d=0$ 时，0 次同调的元素由 0 维图形表示，也就是点。这个情形中的 $\Z$ 表示点的个数。$d=4$ 时，复射影平面中所含的 4 维图形中，没有边界的只有复射影平面自身。这里的 $\Z$ 表示复射影平面的个数。

$d=0,2,4$ 的情形中，$H_d(\CP^2)\cong\Z$ 里与 $1\in\Z$ 对应的元素记作 $u_d$。$u_0$ 是一个点，$u_2$ 是直线所表示的、$\CP^2$ 中所含的图形（因为是在复数意义下考虑的直线，所以实维数是 2），$u_4$ 就是 $\CP^2$ 自身。

好，使用同调的语言，关于相交数的定理 2.2 可以表示如下。
\begin{equation}
n[u_2]\cdot m[u_2] = nm[u_0]. \tag{7}
\end{equation}
这个式子说的是：$n$ 次式所表示的图形（即 $n[u_2]$）与 $m$ 次式所表示的图形（即 $m[u_2]$）的交（用 $\cdot$ 表示）是 $nm$ 个点（即 $nm[u_0]$）。

下面 5 个式子更简单：
\begin{align}
n[u_0]\cdot m[u_0] &= 0, \nonumber \\
n[u_2]\cdot m[u_0] &= 0, \nonumber \\
n[u_4]\cdot m[u_0] &= nm[u_0], \nonumber \\
n[u_4]\cdot m[u_2] &= nm[u_2], \nonumber \\
n[u_4]\cdot m[u_4] &= nm[u_4]. \tag{8}
\end{align}

% ===== PDF p.031 (原书 p.23) =====
例如第 1 个式子说的是：复射影平面上两点的公共部分，除非这两点碰巧重合，否则是空集——这是在说一个理所当然的事实。$n=m=1$ 时的第 3 个式子说的是：复射影平面自身与一点 $P$ 的公共部分就是一点 $P$。

为了把（7）和（8）看得更清楚一些，我们用“Poincaré 对偶”把“同调”改写成“上同调”。Poincaré 对偶是
\begin{equation}
[u_d]^* = x^{(4-d)} \tag{9}
\end{equation}
这样一个映射 $*$，它把 $d$ 维图形所表示的同调的元素 $u_d$，对应到 $4-d$ 次的上同调的元素 $x^{(4-d)}$。

\begin{remark}{3.1}
一般地，对空间 $X$ 可以确定 $d$ 次上同调群 $H^d(X)$ 这样的东西；当 $X$ 是 $n$ 维紧的、有定向的流形时，
\[
H_d(X) \cong H^{n-d}(X)
\]
这样的同构，即 Poincaré 同型，成立。（9）表示的就是这个同构。（复射影平面的情形 $n=4$。）
\end{remark}

那么，用 Poincaré 同型 $*$ 把（7）和（8）搬到上同调一侧，就得到下面的式子。
\begin{equation}
x^{(d)} \cup x^{(d')} =
\begin{cases}
x^{(d+d')}, & \text{当 $d+d' \le 4$ 时}, \\
0, & \text{当 $d+d' \ge 5$ 时}.
\end{cases}
\tag{10}
\end{equation}
这里 $\cup$ 称为杯积，它是相交“数”$\cdot$\footnote{给“数”加上括号，是因为例如（8）的右边不是数，而是同调的元素。}经 Poincaré 同型 $*$ 搬过去得到的东西。例如 $d=d'=2$ 的情形，$x^{(2)}\cup x^{(2)}$ 是
\[
x^{(2)}\cup x^{(2)} = (u_2\cdot u_2)^* = u_0^* = x^{(4)}
\]
（因为 $u_2^* = x^{(2)}$。）
把 $x^{(0)}$ 记作 $1$，$x^{(2)}$ 记作 $x$，$x^{(4)}$ 记作 $x^2$，则
% ===== PDF p.032 (原书 p.24) =====
（10）意味着：在复射影平面的上同调上以 $\cup$ 考虑乘法，得到的是环 $\Z[x]/(x^3)$。这里 $(x^3)$ 指多项式环 $\Z[x]$ 中由 $x^3$ 生成的理想。也就是说，$\Z[x]/(x^3)$ 就是在整数系数一元多项式全体 $\Z[x]$ 上，再加上关系式 $x^3=0$ 得到的东西。总结起来，
\[
H(\CP^2) \cong \Z[x]/(x^3).
\]
这可以推广到维数更高的复射影空间（$n$ 维复射影空间记作 $\CP^n$），有
\begin{equation}
H(\CP^n) \cong \Z[x]/(x^{n+1}) \tag{11}
\end{equation}
成立。下一节想说明的是：当它变成“量子上同调环”时，
\begin{equation}
HQ(\CP^n) \cong \Z[q,x]/(x^{n+1}-q) \tag{12}
\end{equation}
就变成这样。

\section{复射影平面的量子上同调环}

这一节就 $n=2$ 的情形，也就是复射影平面的情形来说明（12）。这个情形下，（12）与（11）的不同之处在于下面两个等式。（请注意，（12）意味着 $x^3=q$ 这个等式成立。）
\begin{equation}
\begin{aligned}
x^1 \cup_q x^2 &= q,\\
x^2 \cup_q x^2 &= qx.
\end{aligned}
\tag{13}
\end{equation}
$\cup_q$ 称为杯积 $\cup$ 的量子形变，或者叫作量子杯积。$q$ 通常被看作新的不定元（也就是形式参数），但（至少在这个情形下）也可以具体代入复数来考虑。代入 $q=0$ 来考虑，（13）就变成 $x^1\cup x^2 = x^2\cup x^2 = 0$，与（11）一致。请这样理解：普通杯积（大体上就是相交数）对应 $q=0$；把 $q$ 换成非零的数，它就“量子形变”为量子杯积。

沿着前一节的最后一步反推回去，就可以把（13）改写成同调的语言。于是这两个式子就是
\begin{equation}
u_2\cdot_q u_0 = qu_4 \tag{14}
\end{equation}
\begin{equation}
u_0\cdot_q u_0 = qu_2 \tag{15}
\end{equation}

% ===== PDF p.033 (原书 p.25) =====
（14）、（15）就是这样。这两个式子所表达的几何内容如下。

\begin{theorem*}{4.1}
$\CP^2$ 的任意两点，都恰好有一条直线通过它们。
\end{theorem*}

定理 4.1 本身是显然的吧。下面说明它为什么会变成（14）、（15）。

设 $\cP,\cQ$ 是 $\CP^2$ 中所含的图形。不带 $q$ 的普通相交“数”，也就是 $\cdot$，大致就是交 $\cP\cap\cQ$。把 $\cdot$ 量子化，得到的 $\cdot_q$ 会是“$q$ 的幂级数”，可以写成
\begin{equation}
\cP\cdot_q\cQ = \sum q^d(\cP\cap_d\cQ) \tag{16}
\end{equation}
这里
\begin{equation}
\cP\cap_d\cQ = \text{与 $\cP$ 以及 $\cQ$ 都相交的 $d$ 次有理曲线的并集}. \tag{17}
\end{equation}
下面来说明“$d$ 次有理曲线”这个词。在 Gromov--Witten 不变量里，我们考虑从黎曼面出发的全纯映射。要定义量子上同调环，就考虑最简单的黎曼面，也就是黎曼球面。黎曼球面就是在复平面上添上无穷远点得到的东西 $\C\cup\{\infty\}$。

复射影平面是在 2 维复向量空间 $\C^2$ 上添上无穷远直线得到的东西。从黎曼球面到复射影平面的全纯映射
\[
\varphi : \C\cup\{\infty\} \longrightarrow \CP^2
\]
可以用以 $z\in\C\cup\{\infty\}$ 为变量的多项式 $f,g,h$ 表示为
\begin{equation}
\varphi(z) = \left(\frac{f(z)}{h(z)},\ \frac{g(z)}{h(z)}\right) \tag{18}
\end{equation}
的样子\footnote{当 $h(z)=0$ 时，认为 $\varphi(z)$ 位于无穷远直线上。}。（可以这样选取，使得不存在满足 $f(a)=g(a)=h(a)=0$ 的 $a$。这样一来，（18）的表示除了把 $f,g,h$ 一起乘上同一个常数外就唯一确定。）

把 $f,g,h$ 看作多项式时，次数中最大的那个数称为 $\varphi$ 的次数；次数为 $d$ 的 $\varphi$ 的像，称为 $d$ 次有理曲线。有理曲线中的“有理”，指的是它是黎曼球面的像。

例如，来看看 $d=0$ 的情形。这个情形 $f,g,h$ 是 0 次多项式，
% ===== PDF p.034 (原书 p.26) =====
也就是常数。所以 $\varphi$ 的像是一点，即 $\varphi$ 是常值映射 $\varphi(z)=c\in\CP^2$。这时，若 $\varphi$ 的像与 $\cP$、$\cQ$ 两者都相交，那么 $c$ 必须同时含于 $\cP$ 与 $\cQ$。也就是说，（17）变成
\begin{equation}
\cP\cap_0\cQ = \cP\cap\cQ \tag{19}
\end{equation}
由此可知，（16）在 $q=0$ 时就是普通的相交“数”。

其次考虑 $d=1$ 的情形。这时 $f,g,h$ 是 1 次式。设
\[
h(z) = a(z-b).
\]
对定义域 $\C\cup\{\infty\}$ 作变量变换，使 $b$ 变成无穷远点 $\infty$\footnote{为此，例如作一次分式变换 $z=b(w+1)/w$ 即可。作为例子，取 $b=1$、$\varphi(z)=(z/(z-1),(z+1)/(z-1))$ 的情形，代入 $z=(w+1)/w$ 计算，就得到 $\varphi(w)=(w+1,2w+1)$。}。于是（18）的分母变成常数，不再需要，$\varphi$ 就可以表示成
\[
\varphi(z) = (az+b,\ cz+d).
\]
也就是说，1 次有理曲线的像就是直线\footnote{这里稍微说了个谎。也就是说，无穷远直线不会变成这样。这是因为在（18）那里说了谎。}。

既然明白了这一点，就来考虑（15）。$u_0$ 是一点。所以考虑 $\cP,\cQ$ 分别是（$\CP^2$ 的）一点 $P$、$Q$ 的情形。那么，（16）的 $d=1$ 的情形，就是通过 $P$ 和 $Q$ 的直线的并集。定理 4.1 说的是：这样的直线恰好有一条，即
\[
u_0\cap_1 u_0 = \text{直线 1 条}
\]
这正是上述几何事实的含义。表示直线的 2 次同调类曾是 $u_2$，所以
\[
u_0\cdot_q u_0 = qu_2 + q^2\ \text{以上的项}
\]
就成了这样。其实在射影平面的量子上同调环的情形，可以验证 $q^2$ 以上出现的项、即 2 次以上有理曲线的贡献为零，于是（15）得证。

接下来考虑（14）。$u_0$ 是一点，$u_2$ 是直线，所以 $\cP$ 是直线，$\cQ$ 是一点 $Q$。那么 $u_0\cap_1 u_2$ 就是与直线 $\cP$ 相交、并且通过 $Q$ 的直线的全体所成的并集。

在 $\cP$ 上取一点 $P$，通过它与 $Q$ 的直线 $PQ$ 恰好有一条（定理 4.1）。让 $P$ 动起来取并集，
\[
\bigcup_{P\in\cP}\overline{PQ}
\]

% ===== PDF p.035 (原书 p.27) =====
这个并集是 2 维复向量空间 $\C^2$ 全体，或者考虑到无穷远点，就是 $\CP^2$ 全体。表示 $\CP^2$ 全体的正是 $u_4$。也就是说，
\[
u_0\cdot_1 u_2 = u_4,
\]
由此就得到了（14）。

\section{计数几何}

前一节考虑了杯积的量子化，也就是 $\cdot_q$。那时对同调群的两个元素 $u,u'$，$u\cdot_q u'$ 仍然是同调群的元素。要得到数而不是同调群的元素，就再取一个同调群的元素 $u''$，考虑
\begin{equation}
q_d(u,u',u'') = (u\cdot_q u')\cdot u'' \tag{20}
\end{equation}
（右边的 $\cdot$ 只是相交数。）回顾第 IV 节的计算，在射影平面的情形可以直接算出
\begin{equation}
q_1(u_a,u_b,u_c) =
\begin{cases}
1, & \text{当 $a+b+c=4$ 时}, \\
0, & \text{其他情形},
\end{cases}
\tag{21}
\end{equation}
结果就是这样\footnote{$u_a$ 是 $a$ 维的图形，$u_0$ 是点，$u_2$ 是直线，$u_4$ 是射影平面全体。}。另一方面，$q_d(u,u',u'')$ 有下面这样的自然的几何意义。
\begin{equation}
q_d(u,u',u'') = \text{“通过 $u,u',u''$ 的 $d$ 次有理曲线的个数”} \tag{22}
\end{equation}
右边只在满足条件的有理曲线为有限个时才有意义，否则取 0。

不只 3 个，还可以考虑更多同调群的元素 $u(1),\cdots,u(k)$。也就是说，
\begin{equation}
q_d(u(1),\cdots,u(k)) = \text{“通过 $u(1),\cdots,u(k)$ 的 $d$ 次有理曲线的个数”} \tag{23}
\end{equation}
这个量 $q_d(u(1),\cdots,u(k))$ 在 $u(i)$ 的置换下不变，而且对 $u(i)$
% ===== PDF p.036 (原书 p.28) =====
是线性的。并且 $u(i)$ 中若有 4 维的（即 $\CP^2$）就变成 0。由此可知，只要考虑 $u(i)$ 是点或直线的情形就够了。由右边 $d$ 次有理曲线的个数为有限个这个条件，只有
\[
q_d(\underbrace{u_0,\cdots,u_0}_{3d-1},\underbrace{u_2,\cdots,u_2}_{1})
\]
是不为 0 的量。另一方面，
\[
q_d(\underbrace{u_0,\cdots,u_0}_{3d-1},\underbrace{u_2,\cdots,u_2}_{1})
= d^1 q_d(\underbrace{u_0,\cdots,u_0}_{3d-1})
\]
% [Tutor Note: 原文此处印作 $d^1$。按除子公理，插入一个直线条件相当于乘上该曲线的次数 $d$；一般地，插入 $l$ 个这样的条件就给出因子 $d^l$。此处是 $l=1$ 的情形，原文照录。]
上式可以直接验证，所以
\begin{equation}
n_d = q_d(\underbrace{u_0,\cdots,u_0}_{3d-1}) \tag{24}
\end{equation}
算出来就足够了。这个数 $n_d$ 的几何意义是
\[
n_d = \text{通过给定的 $3d-1$ 个点的 $d$ 次曲线的个数}
\]
\footnote{不过，“数”这个说法很微妙。$d$ 大时，必然会出现相当于“重根”“数无穷远处的根”的问题，必须好好决定“数”是什么意思。“Schubert 的计数几何的严格基础”这个 Hilbert 第 15 问题，在同样的状况下，也关系到“数”是什么意思这个问题。}。$d=1$ 时这个数是 1，这就是定理 4.1。

求 $n_d$ 这样的数的问题称为“计数几何”（enumerative geometry），在代数几何里是古典的问题。

由这个数决定的母函数满足（最初只提了名字的）WDVV 方程\footnote{这是把量子上同调环的乘法满足结合律这一事实稍微一般化得到的东西。}，利用这一点，可以证明下面的递推式。
\begin{equation}
n_d = \sum_{a+b=d} n_a n_b
\left(a^2b^2\,{}_{3d-4}C_{3a-2} - a^3b\,{}_{3d-4}C_{3a-1}\right). \tag{25}
\end{equation}
（（25）由 Kontsevich 从 WDVV 方程导出。射影平面情形的 WDVV 方程由 Ruan 和 Tian 证明。）更详细的内容，例如参见 [3] 第 6.3 节。

\section{Calabi--Yau 流形}

镜像对称中（或者说超弦理论中）Calabi--Yau 流形起着重要的作用，最后想稍作说明。

上同调的元素带有次数。普通的杯积中，$d$ 次上同调的元素与 $d'$ 次上同调的元素之积是 $d+d'$ 次。这对应于一个几何事实：在 $n$ 维空间中，$m$ 维图形与 $m'$ 维图形的交“通常”是 $m+m'-n$ 维图形。

然而在（13）中，2 次的元素 $x$ 与 4 次的元素 $x^2$ 之积变成了 0 次的元素 $q$（$q$ 是形变的参数，所以是像数一样的东西）。

% ===== PDF p.037 (原书 p.29) =====
因为 $2+4\ne 0$，所以上同调的次数在量子杯积下不被保持。如果要求量子杯积也保持上同调次数，这样的空间就只有 Calabi--Yau 流形。

说明射影平面 $\CP^2$ 的情形下，上同调的次数为什么会在量子杯积下不被保持。用 $\mathcal{M}_d(\CP^2)$ 表示射影平面的 $d$ 次有理曲线的全体所成的集合。这时，
\begin{equation}
\dim \mathcal{M}_d(\CP^2) = 6d-2. \tag{26}
\end{equation}
（26）可以如下证明。射影平面的 $d$ 次有理曲线，是由 $d$ 次多项式 $f,g,h$ 按（18）那样给出映射 $\varphi$，再取它的像得到的。$f,g,h$ 由系数决定，而 $d$ 次多项式的系数有 $d+1$ 个。也就是说，自由度为复 $3d+3$ 维。

但是，把 $f,g,h$ 一起乘上同一个常数，得到的 $\varphi$ 并不改变。也就是说，$\varphi$ 的自由度是复 $3d+2$ 维。

最后，把映射 $\varphi$ 用定义域 $\C\cup\{\infty\}$ 的复解析坐标变换换掉，有理曲线也不改变。黎曼球面 $\C\cup\{\infty\}$ 的复解析坐标变换，由一次分式变换 $z\mapsto (az+b)/(cz+d)$ 表示。它由 4 个复数 $a,b,c,d$ 决定；但把 $a,b,c,d$ 一起乘上同一个常数仍是同一个一次分式变换，所以一次分式变换的自由度为复 3 维。

把这些减掉，射影平面的 $d$ 次有理曲线的自由度就是 $2(3d-1)$，于是得到（26）。

那么，（26）值得注意的一点是：右边随 $d$ 变化。回想 $d=0$ 的情形就是普通杯积，就能明白，这意味着量子杯积不保持次数。

那么，把 $\CP^2$ 换成一般的 $X$，来考虑一下 $\dim \mathcal{M}_d(X)$ 不依赖于 $d$
% ===== PDF p.038 (原书 p.30) =====
的条件吧。

这里只考虑 4 维复射影空间 $\CP^4$ 中的（复）超曲面。首先
\begin{equation}
\dim \mathcal{M}_d(\CP^4) = 10d+2 \tag{27}
\end{equation}
这是因为，$d$ 次的映射
\[
\varphi : \C\cup\{\infty\} \longrightarrow \CP^4
\]
可以表示为
\begin{equation}
\varphi = \left(\frac{h_1(z)}{h_0(z)},\cdots,\frac{h_4(z)}{h_0(z)}\right) \tag{28}
\end{equation}
（$h_0,h_1,\cdots,h_4$ 是 $d$ 次式），与（26）同样可以确认。

设 $\CP^4$ 中的（复）超曲面 $X_m$ 由 4 变量 $m$ 次多项式，例如
\begin{equation}
x^m + y^m + u^m + v^m = 1 \tag{29}
\end{equation}
所表示；来计算 $\mathcal{M}_d(X_m)$ 的维数。因为 $X_m\subset\CP^4$，所以 $\varphi:\C\cup\{\infty\}\to X_m$ 仍然像（28）那样表示。$\varphi$ 的像含于 $X_m$ 的条件是
\[
\left(\frac{h_1(z)}{h_0(z)}\right)^m + \cdots + \left(\frac{h_4(z)}{h_0(z)}\right)^m = 1
\]
是 $z$ 的恒等式。去掉分母就得到
\begin{equation}
(h_1(z))^m + \cdots + (h_4(z))^m = (h_0(z))^m. \tag{30}
\end{equation}
那么，$h_1,\cdots,h_4,h_0$ 是 $d$ 次多项式，它们的系数（有 $5(d+1)$ 个）决定了 $\varphi$。

（30）是关于 $z$ 的 $md$ 次式。它的系数有 $md+1$ 个，这些系数由 $h_1,\cdots,h_4,h_0$ 的系数表示。因此，（30）就在 $h_1,\cdots,h_4,h_0$ 的系数之间给出了 $md+1$ 个条件式。\footnote{需要确认这些条件式全都独立。}
% [Tutor Note: 原文说（30）给出 $md+1$ 个条件式。每个系数取复数值，所以实的条件数是 $2(md+1)$——下文的减法正由此而来。]
因此，$\mathcal{M}_d(X_m)$ 的维数比 $\mathcal{M}_d(\CP^4)$ 的维数少 $2(md+1)$，结果
\[
\dim \mathcal{M}_d(X_m) = (10-2m)d
\]
就得到了。它不依赖于 $d$ 的情形，是 $m=5$ 的时候，也就是 5 次超曲面（quintic）的情形。总结起来，

\begin{theorem*}{6.1}
使 $\CP^4$ 的超曲面的量子杯积保持次数的，是 5 次超曲面的情形。
\end{theorem*}

% ===== PDF p.039 (原书 p.31) =====
第 1 章中已经说明：镜像对称研究开始受到数学家关注，正是通过 5 次超曲面上有理曲线的计数。5 次超曲面在镜像对称研究中之所以重要，理由之一就是定理 6.1。

那么，5 次超曲面是 Calabi--Yau 流形的典型例子。Calabi--Yau 流形是 Ricci 曲率为 0 的 Kähler 流形。下面稍微讲一下定理 6.1 的结论与 Ricci 曲率是怎样联系起来的。

本节对定理 6.1 的说明，是数出表示有理曲线的多项式映射的系数个数、以及它应当满足的方程式的个数，再由两者之差计算维数。这相当于用（有限变量的）多项式来表示 $\mathcal{M}_d(X_m)$。
% [Tutor Note: 原文此处印作「有現変数」，从上下文看疑为「有限変数」（有限个变量）之误，译文按文意处理。]

若用更微分几何的做法，就用微分方程式来表示 $\mathcal{M}_d(X_m)$。这个情形是（非线性）Cauchy--Riemann 方程式。把这个未知变量的个数（实际上是无穷大，因为是映射空间）减去方程式的个数（这也是“某个映射为 0”的方程式，所以也是无穷大），所得的东西称为“差”\footnote{它会是有限的数。}。只要能把这个概念正当化，就可以用它计算 $\mathcal{M}_d(X_m)$ 的维数。做这种正当化的，就是 Atiyah--Singer 指数定理。Atiyah--Singer 指数定理用与空间的曲率有关的量（示性类）来计算“无穷大减无穷大”。考虑我们的情形会出现哪种曲率，就能看出出现的是 Ricci 曲率。

再补充一点。Atiyah--Singer 指数定理中出现的，与其说是曲率本身，不如说是由曲率决定的示性类所表示的上同调类。这个上同调类由（近）复结构决定，不依赖于度量。$\mathcal{M}_d(X_m)$ 的维数计算中出现的是 $X_m$ 的第一 Chern 类。Ricci 曲率为 0 时，第一 Chern 类为 0。第一 Chern 类为 0 的 Kähler 流形上，必定存在 Ricci 曲率为 0 的 Kähler 度量——这就是 Calabi 猜想，证明它的是 Yau。之所以起名叫 Calabi--Yau 流形，就是因为这个。

最后举几篇有“自卖自夸”之嫌的引用文献。

% ===== PDF p.040 (原书 p.32) =====
\begin{reflist}
\item[{[1]}] 深谷賢治，『これからの幾何学』，日本評論社，1998 年．
\item[{[2]}] \texttt{http://www.math.kyoto-u.ac.jp/\~{}fukaya/} に置いてある，1997 年の静岡大学の講義録（横山美佐子さんによる．）
\item[{[3]}] 深谷賢治，『シンプレクティック幾何学』，岩波書店，1999 年．
\item[{[4]}] 中島啓（編），『微分幾何学の最先端』，培風館，2005 年．
\end{reflist}

% END OF CHAPTER 2（PDF p.026--p.040 / 原书 p.18--p.32，共 15 页）
